\chapter{Ordering at a Cafe in Japanese}

% ================= PAIR 1 =================
\begin{convobox}{1}
    \noindent
    \jpkana{みなさん}{Minasan} \jpkana{こんにちは、}{konnichiwa,} 
    \jpkana{たなかです。}{Tanaka desu.}

    \vspace{1.5em}
    \noindent{\Large \color{articolor} \textbf{Arti:} Halo semuanya, saya Tanaka.}
\end{convobox}

\begin{convobox}{2}
    \noindent
    \jpword{きょう}{今日}{Kyou} \jpkana{は、}{wa,} 
    \jpkana{カフェ}{kafe} \jpkana{で}{de} 
    \jpword{ちゅうもん}{注文}{chuumon} \jpkana{するときの}{suru toki no} 
    \jpword{にほんご}{日本語}{Nihongo} \jpkana{を}{o} 
    \jpword{いっしょ}{一緒}{issho} \jpkana{に}{ni} 
    \jpword{れんしゅう}{練習}{renshuu} \jpkana{しましょう。}{shimashou.}

    \vspace{1.5em}
    \noindent{\Large \color{articolor} \textbf{Arti:} Hari ini, mari kita berlatih bahasa Jepang untuk memesan di kafe bersama-sama.}
\end{convobox}

\begin{lessonbox}{✨ Selamat Datang di Kafe Jepang! (Tips Keluarga)}
    Halo Ayah, Bunda, dan Adik-adik! Kafe di Jepang sangat nyaman dan ramah anak. Selain kopi, biasanya selalu ada menu untuk si kecil seperti \textbf{オレンジジュース (Orenji Juusu - Jus Jeruk)} atau \textbf{ミルク (Miruku - Susu sapi)}. Jangan ragu untuk mampir ke kafe saat kalian butuh istirahat sejenak setelah lelah berjalan-jalan!
\end{lessonbox}

% ================= PAIR 2 =================
\begin{convobox}{3}
    \noindent
    \jpkana{カフェ}{Kafe} \jpkana{で}{de} 
    \jpword{ちゅうもん}{注文}{chuumon} \jpkana{するのって、}{suru no tte,} 
    \jpkana{ちょっと}{chotto} 
    \jpword{きんちょう}{緊張}{kinchou} \jpkana{しませんか？}{shimasen ka?}

    \vspace{1.5em}
    \noindent{\Large \color{articolor} \textbf{Arti:} Memesan di kafe itu, rasanya sedikit bikin gugup kan?}
\end{convobox}

\begin{convobox}{4}
    \noindent
    \jpword{じぶん}{自分}{Jibun} \jpkana{の}{no} 
    \jpword{ちゅうもん}{注文}{chuumon} \jpkana{したい}{shitai} \jpword{もの}{物}{mono} \jpkana{が、}{ga,} 
    \jpkana{うまく}{umaku} \jpword{つた}{伝}{tsuta}\jpkana{わらなかった}{waranakatta} 
    \jpword{けいけん}{経験}{keiken} \jpkana{のある}{no aru} 
    \jpword{かた}{方}{kata} \jpkana{もいるかもしれません。}{mo iru kamo shiremasen.}

    \vspace{1.5em}
    \noindent{\Large \color{articolor} \textbf{Arti:} Mungkin ada juga dari kalian yang punya pengalaman (di mana) pesanan yang kalian inginkan tidak tersampaikan dengan baik.}
\end{convobox}

\begin{lessonbox}{💡 Jangan Gugup, Gunakan Bahasa Tubuh!}
    Kalau kamu takut pesananmu salah, jangan panik! Hampir semua kafe di Jepang punya menu bergambar atau etalase kaca berisi kue. Kamu cukup menunjuk gambar atau kuenya dan bilang: \textbf{「これ、お願いします」 (Kore, onegaishimasu - Tolong yang INI)}. Kasir pasti akan langsung mengerti!
\end{lessonbox}

% ================= PAIR 3 =================
\begin{convobox}{5}
    \noindent
    \jpword{てんいん}{店員}{Tenin}\jpkana{さんは}{san wa} 
    \jpword{はやぐち}{早口}{hayaguchi} \jpkana{だし、}{dashi,} 
    \jpword{けいご}{敬語}{keigo} \jpkana{を}{o} 
    \jpword{つか}{使}{tsuka}\jpkana{うので、}{u node,}

    \vspace{1.5em}
    \noindent{\Large \color{articolor} \textbf{Arti:} Kasir (pelayan toko) itu bicaranya cepat, dan karena mereka menggunakan bahasa super sopan (Keigo),}
\end{convobox}

\begin{convobox}{6}
    \noindent
    \jpword{しつもん}{質問}{Shitsumon} \jpkana{を}{o} 
    \jpword{き}{聞}{ki}\jpkana{き}{ki}\jpword{と}{取}{to}\jpkana{るのが}{ru no ga} 
    \jpword{むずか}{難}{muzuka}\jpkana{しいと}{shii to} 
    \jpword{おも}{思}{omo}\jpkana{います。}{imasu.}

    \vspace{1.5em}
    \noindent{\Large \color{articolor} \textbf{Arti:} aku pikir akan sulit untuk menangkap (mendengar) pertanyaan mereka.}
\end{convobox}

\begin{lessonbox}{🛡️ Tameng Anti-Bingung: Dengarkan Kata Kuncinya!}
    Bahasa sopan kasir Jepang (\textit{Keigo}) memang sangat panjang dan cepat. Rahasianya: \textbf{Jangan coba menerjemahkan setiap kata!} Dengarkan saja kata kuncinya. Kalau terdengar kata \textbf{"Hotto"} atau \textbf{"Aisu"}, berarti mereka tanya suhu minuman. Fokus pada kata kunci saja ya!
\end{lessonbox}

% ================= PAIR 4 =================
\begin{convobox}{7}
    \noindent
    \jpkana{そこで}{Soko de} 
    \jpword{きょう}{今日}{kyou} \jpkana{は、}{wa,} 
    \jpkana{カフェ}{kafe} \jpkana{で}{de} \jpkana{よく}{yoku} 
    \jpword{き}{聞}{ki}\jpkana{かれる}{kareru} 
    \jpword{しつもん}{質問}{shitsumon} \jpkana{と}{to} 
    \jpkana{その}{sono} \jpword{こた}{答}{kota}\jpkana{え}{e}\jpword{かた}{方}{kata} \jpkana{を}{o} 
    \jpword{ようい}{用意}{youi} \jpkana{しました。}{shimashita.}

    \vspace{1.5em}
    \noindent{\Large \color{articolor} \textbf{Arti:} Oleh karena itu hari ini, aku menyiapkan pertanyaan yang sering ditanyakan di kafe beserta cara menjawabnya.}
\end{convobox}

\begin{convobox}{8}
    \noindent
    \jpkana{この}{Kono} \jpword{どうが}{動画}{douga} \jpkana{を}{o} 
    \jpword{み}{観}{mi}\jpkana{て}{te} 
    \jpword{れんしゅう}{練習}{renshuu} \jpkana{すれば、}{sureba,} 
    \jpword{じしん}{自信}{jishin} \jpkana{を}{o} 
    \jpword{も}{持}{mo}\jpkana{って}{tte} 
    \jpword{ちゅうもん}{注文}{chuumon} \jpkana{できるように}{dekiru you ni} \jpkana{なります。}{narimasu.}

    \vspace{1.5em}
    \noindent{\Large \color{articolor} \textbf{Arti:} Jika kamu menonton video ini dan berlatih, kamu akan bisa memesan dengan penuh percaya diri.}
\end{convobox}

\begin{lessonbox}{🌟 Percaya Diri Adalah Kunci Utama!}
    Orang Jepang sangat menghargai turis yang mencoba bicara bahasa Jepang, betapapun terbata-batanya. Ajarkan anak-anak untuk selalu tersenyum dan mengucapkan \textbf{「ありがとう」 (Arigatou - Terima kasih)} setelah menerima struk atau minuman. Itu saja sudah membuat kasir sangat senang!
\end{lessonbox}

% ================= PAIR 5 =================
\begin{convobox}{9}
    \noindent
    \jpkana{それでは}{Sore dewa} 
    \jpword{さっそく}{早速}{sassoku} 
    \jpword{い}{行}{i}\jpkana{ってみましょう！}{tte mimashou!}

    \vspace{1.5em}
    \noindent{\Large \color{articolor} \textbf{Arti:} Kalau begitu, ayo langsung kita mulai!}
\end{convobox}

\vspace{1em}
\begin{tcolorbox}[colback=boxframe, colframe=boxframe, rounded corners, boxrule=0pt, arc=3mm, left=2mm, top=2mm, bottom=2mm]
    \Large \bfseries \color{white} ☕ 1. Conversation (Percakapan Dasar)
\end{tcolorbox}
\vspace{1em}

\begin{convobox}{10}
    \noindent
    \jpword{まず}{先ず}{Mazu} \jpword{いちど}{一度}{ichido} 
    \jpkana{カフェ}{kafe} \jpkana{での}{de no} 
    \jpword{かいわ}{会話}{kaiwa} \jpkana{を}{o} 
    \jpword{しぜん}{自然}{shizen} \jpkana{な}{na} 
    \jpkana{スピードで}{supiido de} 
    \jpword{き}{聞}{ki}\jpkana{いてみましょう。}{ite mimashou.}

    \vspace{1.5em}
    \noindent{\Large \color{articolor} \textbf{Arti:} Pertama-tama, mari dengarkan sekali percakapan di kafe dengan kecepatan natural.}
\end{convobox}

\begin{lessonbox}{🤫 Aturan Emas Kafe Jepang: Amankan Kursi Dulu!}
    Sebelum kita memesan di kasir, ada aturan penting di kafe Jepang (seperti Starbucks atau Tully's): \textbf{Cari dan amankan kursi/meja kalian terlebih dahulu!} Kamu bisa meletakkan syal, payung, atau saputangan di atas meja. Setelah meja aman, barulah seluruh keluarga pergi ke kasir untuk memesan.
\end{lessonbox}

% ================= PAIR 6 =================
\begin{convobox}{11}
    \noindent
    \jpkana{いらっしゃいませ、}{Irasshaimase,} 
    \jpword{てんない}{店内}{ten-nai} 
    \jpword{りよう}{利用}{riyou} \jpkana{ですか？}{desu ka?}

    \vspace{1.5em}
    \noindent{\Large \color{articolor} \textbf{Arti:} Selamat datang, apakah dimakan di dalam (Dine-in)?}
\end{convobox}

\begin{convobox}{12}
    \noindent
    \jpkana{いえ、}{Ie,} 
    \jpword{も}{持}{mo}\jpkana{ち}{chi}\jpword{かえ}{帰}{kae}\jpkana{りで}{ri de} 
    \jpkana{お}{o}\jpword{ねが}{願}{nega}\jpkana{いします。}{ishimasu.}

    \vspace{1.5em}
    \noindent{\Large \color{articolor} \textbf{Arti:} Tidak, tolong dibungkus (Take-away).}
\end{convobox}

\begin{lessonbox}{💰 Kenapa Pertanyaan Ini Sangat Penting?}
    Kasir wajib bertanya ini karena \textbf{Pajak di Jepang berbeda!} 
    \begin{itemize}
        \item Bawa pulang (\textit{Mochikaeri}) pajaknya \textbf{8\%}.
        \item Makan di tempat (\textit{Ten-nai}) pajaknya \textbf{10\%}.
    \end{itemize}
    Jadi, pastikan kamu menjawab dengan benar ya agar kasir bisa menghitung total harganya dengan tepat!
\end{lessonbox}

% ================= PAIR 7 =================
\begin{convobox}{13}
    \noindent
    \jpkana{かしこまりました。}{Kashikomarimashita.}

    \vspace{1.5em}
    \noindent{\Large \color{articolor} \textbf{Arti:} Baik, saya mengerti (Sangat Sopan).}
\end{convobox}

\begin{convobox}{14}
    \noindent
    \jpkana{コーヒー}{Koohii} \jpkana{ひとつ}{hitotsu} 
    \jpkana{お}{o}\jpword{ねが}{願}{nega}\jpkana{いします。}{ishimasu.}

    \vspace{1.5em}
    \noindent{\Large \color{articolor} \textbf{Arti:} Tolong pesan satu kopi.}
\end{convobox}

\begin{lessonbox}{🎀 Kata Sopan Andalan Pelayan: Kashikomarimashita}
    Kata \textbf{かしこまりました (Kashikomarimashita)} adalah bahasa yang sangaaaat sopan dari "Saya mengerti" (\textit{Wakarimashita}). Kamu akan mendengar kata ini di hotel, restoran, dan kafe! Sebagai tamu, kamu tidak perlu memakai kata ini, cukup mendengarkan saja.
\end{lessonbox}

% ================= PAIR 8 =================
\begin{convobox}{15}
    \noindent
    \jpkana{ホットか}{Hotto ka} \jpkana{アイス、}{aisu,} 
    \jpkana{どちらに}{dochira ni} \jpkana{なさいますか？}{nasaimasu ka?}

    \vspace{1.5em}
    \noindent{\Large \color{articolor} \textbf{Arti:} Hot (panas) atau Ice (dingin), Anda mau pilih yang mana?}
\end{convobox}

\begin{convobox}{16}
    \noindent
    \jpkana{ホットで}{Hotto de} 
    \jpkana{お}{o}\jpword{ねが}{願}{nega}\jpkana{いします。}{ishimasu.}

    \vspace{1.5em}
    \noindent{\Large \color{articolor} \textbf{Arti:} Tolong yang panas (Hot).}
\end{convobox}

\begin{lessonbox}{🧊 Bahasa Inggris ala Jepang (Japanglish)}
    Hati-hati saat mengucapkan kata bahasa Inggris di Jepang! 
    \begin{itemize}
        \item Jangan bilang "Hot", bilangnya \textbf{ホット (Hotto)}.
        \item Jangan bilang "Ice", bilangnya \textbf{アイス (Aisu)}.
    \end{itemize}
    Aisu juga bisa berarti "Es Krim" di situasi tertentu. Pastikan lafalnya seperti bahasa Jepang ya!
\end{lessonbox}

% ================= PAIR 9 =================
\begin{convobox}{17}
    \noindent
    \jpkana{お}{O}\jpword{さとう}{砂糖}{satou} \jpkana{・}{・} \jpkana{ミルクは}{miruku wa} 
    \jpword{りよう}{利用}{riyou} \jpkana{ですか？}{desu ka?}

    \vspace{1.5em}
    \noindent{\Large \color{articolor} \textbf{Arti:} Apakah menggunakan gula dan susu?}
\end{convobox}

\begin{convobox}{18}
    \noindent
    \jpkana{ミルクだけ}{Miruku dake} 
    \jpkana{お}{o}\jpword{ねが}{願}{nega}\jpkana{いします。}{ishimasu.}

    \vspace{1.5em}
    \noindent{\Large \color{articolor} \textbf{Arti:} Tolong susu saja.}
\end{convobox}

\begin{lessonbox}{☕ Self-Service Area di Kafe}
    Terkadang, kasir tidak menanyakan ini karena ada \textbf{Area Self-Service} (Meja khusus di mana kamu bisa mengambil gula, sedotan, tisu, dan susu cair kecil sendiri). Jika kamu melihat meja penuh toples kecil di dekat kasir, kamu bisa meracik gulamu sendiri di sana!
\end{lessonbox}

% ================= PAIR 10 =================
\begin{convobox}{19}
    \noindent
    \jpkana{お}{O}\jpword{かいけい}{会計}{kaikei}\jpkana{、}{,} 
    \jpword{ごひゃくえん}{500円}{gohyaku-en} 
    \jpkana{でございます。}{de gozaimasu.}

    \vspace{1.5em}
    \noindent{\Large \color{articolor} \textbf{Arti:} Total pembayarannya, 500 yen.}
\end{convobox}

\begin{convobox}{20}
    \noindent
    \jpword{すこ}{少}{suko}\jpkana{し}{shi} 
    \jpword{はや}{速}{haya}\jpkana{かったですか？}{katta desu ka?}

    \vspace{1.5em}
    \noindent{\Large \color{articolor} \textbf{Arti:} Apakah tadi sedikit terlalu cepat?}
\end{convobox}

\begin{lessonbox}{💳 Siapkan Alat Bayar Andalan Keluarga}
    Saat liburan keluarga, membayar dengan uang receh akan memakan waktu. Paling praktis adalah menggunakan kartu IC Card seperti \textbf{Suica} atau \textbf{Pasmo}! Tempelkan kartunya di mesin, berbunyi "Pi!", dan selesai! Sangat praktis untuk beli minuman dan kereta.
\end{lessonbox}

% ================= PAIR 11 =================
\begin{convobox}{21}
    \noindent
    \jpword{いま}{今}{Ima} \jpkana{から}{kara} 
    \jpword{ひと}{一}{hito}\jpkana{つずつ}{tsu zutsu} 
    \jpword{せつめい}{説明}{setsumei} \jpkana{します。}{shimasu.}

    \vspace{1.5em}
    \noindent{\Large \color{articolor} \textbf{Arti:} Mulai sekarang, saya akan menjelaskannya satu per satu.}
\end{convobox}

\vspace{1em}
\begin{tcolorbox}[colback=boxframe, colframe=boxframe, rounded corners, boxrule=0pt, arc=3mm, left=2mm, top=2mm, bottom=2mm]
    \Large \bfseries \color{white} 🍰 2. Phrase 1 (Makan di Tempat / Bawa Pulang)
\end{tcolorbox}
\vspace{1em}

\begin{convobox}{22}
    \noindent
    \jpkana{「いらっしゃいませ、}{"Irasshaimase,} 
    \jpword{てんない}{店内}{ten-nai} 
    \jpword{りよう}{利用}{riyou} \jpkana{ですか？」}{desu ka?"}

    \vspace{1.5em}
    \noindent{\Large \color{articolor} \textbf{Arti:} "Selamat datang, apakah makan/minum di dalam?"}
\end{convobox}

\begin{lessonbox}{🏫 Ten-nai = Di dalam toko}
    Kata \textbf{店内 (Ten-nai)} terdiri dari kanji 'Mise' (Toko) dan 'Uchi' (Dalam). Jadi ini adalah istilah baku untuk menyebut "Makan di tempat" (Dine-in).
\end{lessonbox}

% ================= PAIR 12 =================
\begin{convobox}{23}
    \noindent
    \jpkana{カフェ}{Kafe} \jpkana{で}{de} 
    \jpword{ちゅうもん}{注文}{chuumon} \jpkana{するとき、}{suru toki,} 
    \jpword{いちばん}{一番}{ichiban} \jpword{さいしょ}{最初}{saisho} \jpkana{に}{ni} 

    \vspace{1.5em}
    \noindent{\Large \color{articolor} \textbf{Arti:} Saat memesan di kafe, hal yang paling pertama}
\end{convobox}

\begin{convobox}{24}
    \noindent
    \jpword{き}{聞}{ki}\jpkana{かれる}{kareru} \jpkana{フレーズです。}{fureezu desu.} 
    \jpkana{「}{"}\jpword{てんない}{店内}{Ten-nai} 
    \jpword{りよう}{利用}{riyou} \jpkana{ですか？」}{desu ka?"}

    \vspace{1.5em}
    \noindent{\Large \color{articolor} \textbf{Arti:} didengar (ditanyakan) adalah frasa ini. "Apakah menggunakan area di dalam toko?"}
\end{convobox}

\begin{lessonbox}{🚫 Jangan Makan Sambil Jalan!}
    Jika keluarga memilih "Bawa Pulang" (Mochikaeri), \textbf{Jangan makan/minum sambil berjalan} ya! Di Jepang, makan sambil berjalan di trotoar dianggap kurang sopan. Carilah bangku taman, atau menyingkir ke sisi jalan untuk meminum jus/kopi kalian!
\end{lessonbox}

% ================= PAIR 13 =================
\begin{convobox}{25}
    \noindent
    \jpkana{もしくは}{Moshiku wa}

    \vspace{1.5em}
    \noindent{\Large \color{articolor} \textbf{Arti:} Atau (bisa juga)}
\end{convobox}

\begin{convobox}{26}
    \noindent
    \jpkana{「お}{"O}\jpword{も}{持}{mo}\jpkana{ち}{chi}\jpword{かえ}{帰}{kae}\jpkana{りですか？」}{ri desu ka?"}

    \vspace{1.5em}
    \noindent{\Large \color{articolor} \textbf{Arti:} "Apakah dibungkus bawa pulang?"}
\end{convobox}

\begin{lessonbox}{🛍️ "O-mochi-kaeri" = Bawa pulang}
    Kata ini asalnya dari kata "Motsu" (Membawa) dan "Kaeru" (Pulang). Ingat baik-baik kata \textbf{持ち帰り (Mochikaeri)} karena kamu akan melihat tulisan ini di kafe, restoran sushi, hingga kedai burger!
\end{lessonbox}

% ================= PAIR 14 =================
\begin{convobox}{27}
    \noindent
    \jpword{こた}{答}{Kota}\jpkana{えるときは、}{eru toki wa,} 
    \jpkana{このように}{kono you ni} \jpword{い}{言}{i}\jpkana{います。}{imasu.}

    \vspace{1.5em}
    \noindent{\Large \color{articolor} \textbf{Arti:} Saat menjawab, ucapkanlah seperti ini.}
\end{convobox}

\begin{convobox}{28}
    \noindent
    \jpkana{「}{"}\jpword{てんない}{店内}{Ten-nai} \jpkana{で}{de} 
    \jpkana{お}{o}\jpword{ねが}{願}{nega}\jpkana{いします。」}{ishimasu."}

    \vspace{1.5em}
    \noindent{\Large \color{articolor} \textbf{Arti:} "Tolong di dalam sini (Makan di tempat)."}
\end{convobox}

\begin{lessonbox}{🪄 Kekuatan Partikel "De"}
    Jangan lupa tambahkan partikel \textbf{で (De)} setelah pilihanmu! \textit{Ten-nai DE} (DI dalam), atau \textit{Mochikaeri DE} (DENGAN cara dibawa pulang). Lalu tutup dengan jurus \textit{Onegaishimasu} agar terdengar manis dan sopan!
\end{lessonbox}

% ================= PAIR 15 =================
\begin{convobox}{29}
    \noindent
    \jpkana{「お}{"O}\jpword{も}{持}{mo}\jpkana{ち}{chi}\jpword{かえ}{帰}{kae}\jpkana{りで}{ri de} 
    \jpkana{お}{o}\jpword{ねが}{願}{nega}\jpkana{いします。」}{ishimasu."}

    \vspace{1.5em}
    \noindent{\Large \color{articolor} \textbf{Arti:} "Tolong dibungkus (Bawa pulang)."}
\end{convobox}

\vspace{1em}
\begin{tcolorbox}[colback=boxframe, colframe=boxframe, rounded corners, boxrule=0pt, arc=3mm, left=2mm, top=2mm, bottom=2mm]
    \Large \bfseries \color{white} 🍩 3. Phrase 2 (Sebut Pesanan \& Jumlah)
\end{tcolorbox}
\vspace{1em}

\begin{convobox}{30}
    \noindent
    \jpkana{「コーヒー}{"Koohii} \jpkana{ひとつ}{hitotsu} 
    \jpkana{お}{o}\jpword{ねが}{願}{nega}\jpkana{いします。」}{ishimasu."}

    \vspace{1.5em}
    \noindent{\Large \color{articolor} \textbf{Arti:} "Tolong 1 buah kopi."}
\end{convobox}

\begin{lessonbox}{🥤 Cara Memesan Jus untuk Anak}
    Jika anak-anak ingin memesan sendiri, ajarkan mereka menyebut minumannya lalu angka \textit{hitotsu}! Contoh:
    "Ringo juusu (Jus apel) hitotsu, onegaishimasu!" Pasti kasirnya akan merasa sangat gemas!
\end{lessonbox}

% ================= PAIR 16 =================
\begin{convobox}{31}
    \noindent
    \jpword{ちゅうもん}{注文}{Chuumon} \jpkana{するものの}{suru mono no} 
    \jpword{なまえ}{名前}{namae} \jpkana{と}{to} 
    \jpword{かず}{数}{kazu} \jpkana{を}{o} 
    \jpword{つた}{伝}{tsuta}\jpkana{える}{eru} 
    \jpkana{フレーズです。}{fureezu desu.}

    \vspace{1.5em}
    \noindent{\Large \color{articolor} \textbf{Arti:} Ini adalah frasa untuk menyampaikan nama dan jumlah barang yang dipesan.}
\end{convobox}

\begin{convobox}{32}
    \noindent
    \jpword{かず}{数}{Kazu} \jpkana{は、}{wa,} 
    \jpkana{このように}{kono you ni} 
    \jpword{い}{言}{i}\jpkana{います。}{imasu.}

    \vspace{1.5em}
    \noindent{\Large \color{articolor} \textbf{Arti:} Untuk jumlahnya, diucapkan seperti ini.}
\end{convobox}

\begin{lessonbox}{📝 Urutan yang Benar}
    Berbeda dengan bahasa Indonesia "Satu Kopi", di Jepang urutannya dibalik: \textbf{BARANGNYA DULU, baru JUMLAHNYA}. Jadi: Kopi 1 (Koohii hitotsu), Kue 2 (Keeki futatsu). Ingat-ingat rumus ini ya!
\end{lessonbox}

% ================= PAIR 17 =================
\begin{convobox}{33}
    \noindent
    \jpkana{ひとつ}{Hitotsu} (1) / \jpkana{ふたつ}{Futatsu} (2) / \jpkana{みっつ}{Mittsu} (3) / \jpkana{よっつ}{Yottsu} (4) / \jpkana{いつつ}{Itsutsu} (5) / \jpkana{むっつ}{Muttsu} (6)
\end{convobox}

\begin{convobox}{34}
    \noindent
    \jpkana{「ひとつ」と}{"Hitotsu" to} 
    \jpkana{「ふたつ」は}{"futatsu" wa} 
    \jpword{はつおん}{発音}{hatsuon} \jpkana{が}{ga} 
    \jpkana{よく}{yoku} \jpword{に}{似}{ni}\jpkana{ているので、}{te iru node,}

    \vspace{1.5em}
    \noindent{\Large \color{articolor} \textbf{Arti:} Karena pengucapan "Hitotsu (1)" dan "Futatsu (2)" itu sangat mirip,}
\end{convobox}

\begin{lessonbox}{Hitotsu dan Futatsu}
    "Hitotsu" dan "Futatsu" terdengar hampir sama, terutama jika kafe sedang ramai dan berisik. Kadang kasir salah dengar kamu pesan 1 jadi pesan 2! Untuk menghindarinya, \textbf{Selalu angkat jarimu} saat menyebutkan angkanya ya!
\end{lessonbox}

% ================= PAIR 18 =================
\begin{convobox}{35}
    \noindent
    \jpword{ゆび}{指}{Yubi} \jpkana{も}{mo} 
    \jpword{だ}{出}{da}\jpkana{すと}{su to} 
    \jpkana{いいでしょう。}{ii deshou.}

    \vspace{1.5em}
    \noindent{\Large \color{articolor} \textbf{Arti:} ada baiknya jika kamu mengeluarkan (menunjukkan) jari juga.}
\end{convobox}

\begin{convobox}{36}
    \noindent
    \jpword{ちゅうもん}{注文}{Chuumon} \jpkana{するものが}{suru mono ga} 
    \jpkana{いくつか}{ikutsu ka} \jpkana{ある}{aru} 
    \jpword{とき}{時}{toki}\jpkana{は、}{wa,} 
    \jpkana{このように}{kono you ni} \jpword{い}{言}{i}\jpkana{います。}{imasu.}

    \vspace{1.5em}
    \noindent{\Large \color{articolor} \textbf{Arti:} Saat barang yang dipesan ada beberapa, ucapkanlah seperti ini.}
\end{convobox}

\begin{lessonbox}{👨‍👩‍👧‍👦 Pesan untuk Seluruh Keluarga}
    Biasanya Ayah atau Bunda yang akan memesan untuk seluruh anggota keluarga kan? Nah, siap-siap merangkai kata dengan partikel penghubung di bawah ini ya!
\end{lessonbox}

% ================= PAIR 19 =================
\begin{convobox}{37}
    \noindent
    \jpkana{「コーヒー}{"Koohii} \jpkana{ひとつ}{hitotsu} \jpkana{と、}{to,} 
    \jpkana{チーズケーキ}{chiizu keeki} \jpkana{ひとつ}{hitotsu} 
    \jpkana{お}{o}\jpword{ねが}{願}{nega}\jpkana{いします。」}{ishimasu."}

    \vspace{1.5em}
    \noindent{\Large \color{articolor} \textbf{Arti:} "Tolong 1 kopi DAN 1 cheesecake."}
\end{convobox}

\begin{convobox}{38}
    \noindent
    \jpkana{「カフェオレ}{"Kafe ore} \jpkana{ふたつ}{futatsu} \jpkana{と、}{to,} 
    \jpkana{サンドイッチ}{sandoitchi} \jpkana{ひとつ}{hitotsu} 
    \jpkana{お}{o}\jpword{ねが}{願}{nega}\jpkana{いします。」}{ishimasu."}

    \vspace{1.5em}
    \noindent{\Large \color{articolor} \textbf{Arti:} "Tolong 2 cafe au lait (kopi susu) DAN 1 sandwich."}
\end{convobox}

\begin{lessonbox}{🔗 Lem Perekat Kalimat: Partikel "To"}
    Partikel \textbf{と (to)} fungsinya sama seperti kata "Dan". Kamu bisa menyambung sebanyak apapun pesananmu! \textit{A to, B to, C to, D, onegaishimasu!}. Kasir akan dengan sabar mencatat semuanya!
\end{lessonbox}

% ================= PAIR 20 =================
\vspace{1em}
\begin{tcolorbox}[colback=boxframe, colframe=boxframe, rounded corners, boxrule=0pt, arc=3mm, left=2mm, top=2mm, bottom=2mm]
    \Large \bfseries \color{white} 🥤 4. Phrase 3 (Hot or Ice?)
\end{tcolorbox}
\vspace{1em}

\begin{convobox}{39}
    \noindent
    \jpkana{「ホットか}{"Hotto ka} \jpkana{アイス、}{aisu,} 
    \jpkana{どちらに}{dochira ni} \jpkana{なさいますか？」}{nasaimasu ka?"}

    \vspace{1.5em}
    \noindent{\Large \color{articolor} \textbf{Arti:} "Hot atau Ice, Anda ingin yang mana?"}
\end{convobox}

\begin{convobox}{40}
    \noindent
    \jpkana{コーヒーや}{Koohii ya} 
    \jpword{こうちゃ}{紅茶}{koucha} \jpkana{のように、}{no you ni,} 
    \jpkana{ホットか}{hotto ka} \jpkana{アイスを}{aisu o} 
    \jpword{えら}{選}{era}\jpkana{べる}{beru} 

    \vspace{1.5em}
    \noindent{\Large \color{articolor} \textbf{Arti:} Seperti kopi atau teh hitam, (minuman) yang bisa dipilih Hot atau Ice-nya}
\end{convobox}

\begin{lessonbox}{🧊 Es di Jepang Itu BANYAK Banget!}
    Sedikit tips: Kalau kamu memesan minuman dingin (\textbf{Aisu}), kafe Jepang biasanya mengisi gelasmu dengan ES BATU yang saaaaaangat penuh! Kalau kamu atau si kecil tidak mau terlalu dingin, bilang saja: \textbf{「氷 少なめで」 (Koori sukuname de - Esnya sedikit saja ya)}.
\end{lessonbox}

% ================= PAIR 21 =================
\begin{convobox}{41}
    \noindent
    \jpword{の}{飲}{no}\jpkana{み}{mi}\jpword{もの}{物}{mono} \jpkana{を}{o} 
    \jpword{ちゅうもん}{注文}{chuumon} \jpkana{した}{shita} 
    \jpword{とき}{時}{toki}\jpkana{は}{wa} 
    \jpkana{このように}{kono you ni} 
    \jpword{き}{聞}{ki}\jpkana{かれます。}{karemasu.}

    \vspace{1.5em}
    \noindent{\Large \color{articolor} \textbf{Arti:} Saat memesan minuman, kamu akan ditanya seperti ini.}
\end{convobox}

\begin{convobox}{42}
    \noindent
    \jpword{こた}{答}{Kota}\jpkana{えるときは、}{eru toki wa,} 
    \jpkana{こう}{kou} \jpword{い}{言}{i}\jpkana{いましょう。}{imashou.}

    \vspace{1.5em}
    \noindent{\Large \color{articolor} \textbf{Arti:} Saat menjawab, mari ucapkan seperti ini.}
\end{convobox}

\begin{lessonbox}{☕ Pilihan Cerdas Musim Dingin}
    Kalau kamu liburan ke Jepang saat musim dingin (Winter), memesan \textbf{Hotto} (Minuman panas) adalah pilihan terbaik! Gelas kopimu juga bisa berfungsi ganda sebagai penghangat tangan (Kairo) dadakan!
\end{lessonbox}

% ================= PAIR 22 =================
\begin{convobox}{43}
    \noindent
    \jpkana{「ホットで}{"Hotto de} \jpkana{お}{o}\jpword{ねが}{願}{nega}\jpkana{いします。」}{ishimasu."}

    \vspace{1.5em}
    \noindent{\Large \color{articolor} \textbf{Arti:} "Tolong yang panas."}
\end{convobox}

\begin{convobox}{44}
    \noindent
    \jpkana{「アイスで}{"Aisu de} \jpkana{お}{o}\jpword{ねが}{願}{nega}\jpkana{いします。」}{ishimasu."}

    \vspace{1.5em}
    \noindent{\Large \color{articolor} \textbf{Arti:} "Tolong yang es (dingin)."}
\end{convobox}

\begin{lessonbox}{⚠️ Alergi Susu? Peringatan Penting!}
    Kalau ada anggota keluarga yang alergi susu sapi (Lactose Intolerant), pastikan untuk memesan kopi hitam (Burakku Koohii), dan jangan memesan "Kafe Rate" atau "Cafe o-lait". Atau bilang: \textbf{ミルクのアレルギーがあります (Miruku no arerugii ga arimasu - Saya ada alergi susu)}.
\end{lessonbox}

% ================= PAIR 23 =================
\vspace{1em}
\begin{tcolorbox}[colback=boxframe, colframe=boxframe, rounded corners, boxrule=0pt, arc=3mm, left=2mm, top=2mm, bottom=2mm]
    \Large \bfseries \color{white} 🍯 5. Phrase 4 (Gula \& Susu)
\end{tcolorbox}
\vspace{1em}

\begin{convobox}{45}
    \noindent
    \jpkana{「お}{"O}\jpword{さとう}{砂糖}{satou} \jpkana{・}{・} \jpkana{ミルクは}{miruku wa} 
    \jpword{りよう}{利用}{riyou} \jpkana{ですか？」}{desu ka?"}

    \vspace{1.5em}
    \noindent{\Large \color{articolor} \textbf{Arti:} "Apakah membutuhkan gula dan susu?"}
\end{convobox}

\begin{convobox}{46}
    \noindent
    \jpword{こた}{答}{Kota}\jpkana{えるときは、}{eru toki wa,} 
    \jpkana{このように}{kono you ni} \jpword{い}{言}{i}\jpkana{います。}{imasu.}

    \vspace{1.5em}
    \noindent{\Large \color{articolor} \textbf{Arti:} Saat menjawab, ucapkanlah seperti ini.}
\end{convobox}

\begin{lessonbox}{🍯 Apa Itu "Susu" untuk Kopi Hitam?}
    Di Jepang, kalau kamu pesan kopi hitam lalu kasir menawarkan "Miruku" (Susu), ini BUKAN susu cair literan lho! Ini adalah krimer cair kecil di dalam cup plastik mungil yang disebut \textbf{コーヒーフレッシュ (Koohii Furesshu - Coffee Fresh)}. Gratis kok!
\end{lessonbox}

% ================= PAIR 24 =================
\begin{convobox}{47}
    \noindent
    \jpkana{「はい、}{"Hai,} \jpkana{お}{o}\jpword{ねが}{願}{nega}\jpkana{いします。」}{ishimasu."}

    \vspace{1.5em}
    \noindent{\Large \color{articolor} \textbf{Arti:} "Ya, tolong (berikan keduanya)."}
\end{convobox}

\begin{convobox}{48}
    \noindent
    \jpkana{「}{"}\jpword{さとう}{砂糖}{Satou} \jpkana{だけ}{dake} 
    \jpkana{お}{o}\jpword{ねが}{願}{nega}\jpkana{いします。」}{ishimasu."}

    \vspace{1.5em}
    \noindent{\Large \color{articolor} \textbf{Arti:} "Tolong gula saja."}
\end{convobox}

\begin{lessonbox}{💎 Kata Ajaib "Dake" (Hanya)}
    \textbf{だけ (Dake)} berarti "Hanya/Saja". Kalau anakmu hanya suka sirup strawberry tanpa krim di atas minumannya, kamu bisa bilang \textit{"Ichigo shiroppu \textbf{dake} onegaishimasu"}. Sangat berguna untuk pilih-pilih makanan!
\end{lessonbox}

% ================= PAIR 25 =================
\begin{convobox}{49}
    \noindent
    \jpkana{「ミルクだけ}{"Miruku dake} \jpkana{お}{o}\jpword{ねが}{願}{nega}\jpkana{いします。」}{ishimasu."}

    \vspace{1.5em}
    \noindent{\Large \color{articolor} \textbf{Arti:} "Tolong susu saja."}
\end{convobox}

\begin{convobox}{50}
    \noindent
    \jpkana{「いえ、}{"Ie,} \jpword{だいじょうぶ}{大丈夫}{daijoubu} \jpkana{です。」}{desu."}

    \vspace{1.5em}
    \noindent{\Large \color{articolor} \textbf{Arti:} "Tidak, tidak usah."}
\end{convobox}

\begin{lessonbox}{♻️ Mengembalikan Nampan Sendiri}
    \textbf{BUDAYA PENTING!} Di kafe Jepang seperti Starbucks, setelah selesai makan, kita \textbf{TIDAK BOLEH} meninggalkan gelas dan nampan di meja! Carilah tempat bernama \textbf{返却口 (Henkyakuguchi - Tempat Pengembalian)}. Buang sampahmu dan letakkan nampan di sana. Ini bisa jadi pelajaran disiplin yang seru buat anak-anak!
\end{lessonbox}

% ================= PAIR 26 =================
\vspace{1em}
\begin{tcolorbox}[colback=boxframe, colframe=boxframe, rounded corners, boxrule=0pt, arc=3mm, left=2mm, top=2mm, bottom=2mm]
    \Large \bfseries \color{white} 🗣️ 6. Role Play 1 (Simulasi Memesan Kopi)
\end{tcolorbox}
\vspace{1em}

\begin{convobox}{51}
    \noindent
    \jpkana{それでは、}{Sore dewa,} 
    \jpword{いま}{今}{ima} \jpkana{から}{kara} 
    \jpkana{ロールプレイを}{rooru purei o} \jpkana{します。}{shimasu.}

    \vspace{1.5em}
    \noindent{\Large \color{articolor} \textbf{Arti:} Kalau begitu, sekarang kita akan mulai bermain peran (Roleplay).}
\end{convobox}

\begin{convobox}{52}
    \noindent
    \jpword{さき}{先}{Sakki} \jpkana{ほどの}{hodo no} 
    \jpkana{カフェ}{kafe} \jpkana{での}{de no} 
    \jpword{かいわ}{会話}{kaiwa} \jpkana{を}{o} 
    \jpword{もういちど}{もう一度}{mou ichido} 
    \jpkana{なが}{流}{naga}\jpkana{します。}{shimasu.}

    \vspace{1.5em}
    \noindent{\Large \color{articolor} \textbf{Arti:} Aku akan memutar kembali percakapan kafe yang tadi.}
\end{convobox}

\begin{lessonbox}{🎭 Waktunya Akting! Bersiaplah!}
    Ayo semuanya tarik napas panjang! Tegakkan badanmu! Bayangkan kamu sedang berdiri di depan kasir kafe Jepang di Tokyo. Berbicaralah dengan lantang dan senyum yang ramah ya!
\end{lessonbox}

% ================= PAIR 27 =================
\begin{convobox}{53}
    \noindent
    \jpword{こんかい}{今回}{Konkai} \jpkana{は、}{wa,} 
    \jpkana{みなさんは}{minasan wa} 
    \jpkana{お}{o}\jpword{きゃく}{客}{kyaku}\jpkana{さんに}{san ni} 
    \jpkana{なったつもりで}{natta tsumori de} 

    \vspace{1.5em}
    \noindent{\Large \color{articolor} \textbf{Arti:} Kali ini, bayangkanlah kalian menjadi pelanggannya, dan}
\end{convobox}

\begin{convobox}{54}
    \noindent
    \jpword{こえ}{声}{koe} \jpkana{に}{ni} \jpword{だ}{出}{da}\jpkana{して}{shite} 
    \jpword{はな}{話}{hana}\jpkana{してみてください。}{shite mite kudasai.}

    \vspace{1.5em}
    \noindent{\Large \color{articolor} \textbf{Arti:} cobalah berbicara dengan mengeluarkan suara.}
\end{convobox}

\begin{lessonbox}{🎙️ Suara Lantang itu Bagus!}
    Kafe di Jepang biasanya memutar musik BGM (Background Music) dan ada suara mesin espresso yang bising. Jadi pastikan volume suaramu jelas ya agar kasir tidak meminta kita mengulang pesanan.
\end{lessonbox}

% ================= PAIR 28 =================
\begin{convobox}{55}
    \noindent
    \jpkana{【店員】いらっしゃいませ、}{[Tenin] Irasshaimase,} 
    \jpword{てんない}{店内}{ten-nai} 
    \jpword{りよう}{利用}{riyou} \jpkana{ですか？}{desu ka?}

    \vspace{1.5em}
    \noindent{\Large \color{articolor} \textbf{Arti:} [Kasir] Selamat datang, apakah dimakan di dalam?}
\end{convobox}

\begin{convobox}{56}
    \noindent
    \jpkana{【あなた】いえ、}{[Anata] Ie,} 
    \jpword{も}{持}{mo}\jpkana{ち}{chi}\jpword{かえ}{帰}{kae}\jpkana{りで}{ri de} 
    \jpkana{お}{o}\jpword{ねが}{願}{nega}\jpkana{いします。}{ishimasu.}

    \vspace{1.5em}
    \noindent{\Large \color{articolor} \textbf{Arti:} [Kamu] Tidak, tolong dibungkus bawa pulang.}
\end{convobox}

\begin{lessonbox}{🚶‍♂️ Mochikaeri: Cara Menghemat Waktu!}
    Kalau keluarga harus segera mengejar jadwal kereta Shinkansen, opsi \textit{Mochikaeri} (Bawa Pulang) adalah pilihan terbaik! Kopimu akan dimasukkan ke dalam kantong kertas (Kami-bukuro) yang aman agar tidak tumpah.
\end{lessonbox}

% ================= PAIR 29 =================
\begin{convobox}{57}
    \noindent
    \jpkana{【店員】かしこまりました。}{[Tenin] Kashikomarimashita.}

    \vspace{1.5em}
    \noindent{\Large \color{articolor} \textbf{Arti:} [Kasir] Baik.}
\end{convobox}

\begin{convobox}{58}
    \noindent
    \jpkana{【あなた】コーヒー}{[Anata] Koohii} \jpkana{ひとつ}{hitotsu} 
    \jpkana{お}{o}\jpword{ねが}{願}{nega}\jpkana{いします。}{ishimasu.}

    \vspace{1.5em}
    \noindent{\Large \color{articolor} \textbf{Arti:} [Kamu] Tolong 1 buah kopi.}
\end{convobox}

\begin{lessonbox}{☝️ Jangan Lupa Angkat 1 Jari!}
    Sambil bilang "Hitotsu", angkat satu jarimu (jari telunjuk)! Bahasa tubuh ini sangat universal dan langsung dipahami oleh kasir Jepang.
\end{lessonbox}

% ================= PAIR 30 =================
\begin{convobox}{59}
    \noindent
    \jpkana{【店員】ホットか}{[Tenin] Hotto ka} \jpkana{アイス、}{aisu,} 
    \jpkana{どちらに}{dochira ni} \jpkana{なさいますか？}{nasaimasu ka?}

    \vspace{1.5em}
    \noindent{\Large \color{articolor} \textbf{Arti:} [Kasir] Hot atau Ice, mau pilih yang mana?}
\end{convobox}

\begin{convobox}{60}
    \noindent
    \jpkana{【あなた】ホットで}{[Anata] Hotto de} 
    \jpkana{お}{o}\jpword{ねが}{願}{nega}\jpkana{いします。}{ishimasu.}

    \vspace{1.5em}
    \noindent{\Large \color{articolor} \textbf{Arti:} [Kamu] Tolong yang Hot (Panas).}
\end{convobox}

\begin{lessonbox}{❄️ Hati-Hati Cangkir Panas!}
    Jika kamu memesan \textit{Hotto}, biasanya cangkir kertasnya akan sangat panas. Pelayan kadang akan membungkusnya dengan \textbf{スリーブ (Suriibu - Sleeve / Pelindung kardus)}. Jangan biarkan anak-anak yang memegangnya ya!
\end{lessonbox}

% ================= PAIR 31 =================
\begin{convobox}{61}
    \noindent
    \jpkana{【店員】お}{[Tenin] O}\jpword{さとう}{砂糖}{satou} \jpkana{・}{・} \jpkana{ミルクは}{miruku wa} 
    \jpword{りよう}{利用}{riyou} \jpkana{ですか？}{desu ka?}

    \vspace{1.5em}
    \noindent{\Large \color{articolor} \textbf{Arti:} [Kasir] Apakah butuh gula dan susu?}
\end{convobox}

\begin{convobox}{62}
    \noindent
    \jpkana{【あなた】ミルクだけ}{[Anata] Miruku dake} 
    \jpkana{お}{o}\jpword{ねが}{願}{nega}\jpkana{いします。}{ishimasu.}

    \vspace{1.5em}
    \noindent{\Large \color{articolor} \textbf{Arti:} [Kamu] Tolong susu saja.}
\end{convobox}

\begin{lessonbox}{🍭 Meminta Sedotan atau Tisu Basah}
    Kalau kamu butuh sesuatu yang lain, kamu bisa minta ke kasir!
    Sedotan = \textbf{Sutoro}. Tisu basah untuk lap tangan = \textbf{Oshibori}. 
    Bilang saja: \textit{"Sutoro kudasai!"} atau \textit{"Oshibori kudasai!"}
\end{lessonbox}

% ================= PAIR 32 =================
\begin{convobox}{63}
    \noindent
    \jpkana{【店員】お}{[Tenin] O}\jpword{かいけい}{会計}{kaikei}\jpkana{、}{,} 
    \jpword{ごひゃくえん}{500円}{gohyaku-en} \jpkana{でございます。}{de gozaimasu.}

    \vspace{1.5em}
    \noindent{\Large \color{articolor} \textbf{Arti:} [Kasir] Totalnya 500 yen.}
\end{convobox}

\begin{convobox}{64}
    \noindent
    \jpkana{いい}{Ii} \jpword{かん}{感}{kan}\jpkana{じです！}{ji desu!}

    \vspace{1.5em}
    \noindent{\Large \color{articolor} \textbf{Arti:} [Tanaka Sensei] Kelihatan bagus / Kamu melakukannya dengan baik!}
\end{convobox}

\begin{lessonbox}{📱 Saat Kasir Menunjuk Layar Kecil}
    Jika kamu tidak mengerti angka dalam bahasa Jepang (seperti Gohyaku-en), jangan khawatir! Mesin kasir pasti punya \textbf{Layar Kecil yang menghadap ke arah pembeli}. Cukup lihat angka di layar itu untuk tahu berapa yang harus dibayar. Canggih, kan?
\end{lessonbox}

% ================= PAIR 33 =================
\vspace{1em}
\begin{tcolorbox}[colback=boxframe, colframe=boxframe, rounded corners, boxrule=0pt, arc=3mm, left=2mm, top=2mm, bottom=2mm]
    \Large \bfseries \color{white} 🗣️ 7. Role Play 2 (Pesanan Keluarga Besar!)
\end{tcolorbox}
\vspace{1em}

\begin{convobox}{65}
    \noindent
    \jpword{さいご}{最後}{Saigo} \jpkana{に、}{ni,} 
    \jpword{もういちど}{もう一度}{mou ichido} 
    \jpkana{ロールプレイを}{rooru purei o} \jpkana{します。}{shimasu.}

    \vspace{1.5em}
    \noindent{\Large \color{articolor} \textbf{Arti:} Terakhir, mari kita lakukan roleplay sekali lagi.}
\end{convobox}

\begin{convobox}{66}
    \noindent
    \jpword{こんど}{今度}{Kondo} \jpkana{は、}{wa,} 
    \jpkana{イラストを}{irasuto o} \jpword{み}{見}{mi}\jpkana{て}{te} 
    \jpword{ちゅうもん}{注文}{chuumon} \jpkana{してください。}{shite kudasai.}

    \vspace{1.5em}
    \noindent{\Large \color{articolor} \textbf{Arti:} Kali ini, silakan memesan dengan melihat ilustrasi gambarnya.}
\end{convobox}

\begin{lessonbox}{📸 Menu Jepang Sangat Membantu Visual!}
    Ilustrasi atau foto di menu Jepang 99\% sama dengan bentuk aslinya! Jadi kalau kamu lihat kue cantik dengan stroberi, yang datang nanti persis seperti itu. Jangan ragu pesan dari gambar!
\end{lessonbox}

% ================= PAIR 34 =================
\begin{convobox}{67}
    \noindent
    \jpword{じっさい}{実際}{Jissai} \jpkana{に}{ni} 
    \jpkana{カフェにいるのを}{kafe ni iru no o} 
    \jpkana{イメージして}{imeeji shite} 
    \jpword{はな}{話}{hana}\jpkana{してみてくださいね。}{shite mite kudasai ne.}

    \vspace{1.5em}
    \noindent{\Large \color{articolor} \textbf{Arti:} Cobalah berbicara sambil mengimajinasikan bahwa kamu benar-benar sedang berada di kafe (sungguhan).}
\end{convobox}

\begin{convobox}{68}
    \noindent
    \jpkana{【店員】いらっしゃいませ、}{[Tenin] Irasshaimase,} 
    \jpword{てんない}{店内}{ten-nai} 
    \jpword{りよう}{利用}{riyou} \jpkana{ですか？}{desu ka?}

    \vspace{1.5em}
    \noindent{\Large \color{articolor} \textbf{Arti:} [Kasir] Selamat datang, dimakan di sini?}
\end{convobox}

\begin{lessonbox}{🎒 Anak-anak Berbaris di Belakang}
    Tips keluarga: Saat memesan di kasir yang sempit, biarkan satu orang tua saja yang maju berbicara ke kasir (membawa catatan pesanan), sedangkan anak-anak dan orang tua satunya bisa menunggu atau mencari kursi terlebih dahulu!
\end{lessonbox}

% ================= PAIR 35 =================
\begin{convobox}{69}
    \noindent
    \jpkana{【あなた】はい、}{[Anata] Hai,} 
    \jpword{てんない}{店内}{ten-nai} \jpkana{で}{de} 
    \jpkana{お}{o}\jpword{ねが}{願}{nega}\jpkana{いします。}{ishimasu.}

    \vspace{1.5em}
    \noindent{\Large \color{articolor} \textbf{Arti:} [Kamu] Ya, tolong dimakan di sini.}
\end{convobox}

\begin{convobox}{70}
    \noindent
    \jpkana{【店員】かしこまりました。}{[Tenin] Kashikomarimashita.}

    \vspace{1.5em}
    \noindent{\Large \color{articolor} \textbf{Arti:} [Kasir] Baik.}
\end{convobox}

\begin{lessonbox}{🍵 Air Putih Gratis (O-hiya)}
    Saat kamu makan di tempat (\textit{Ten-nai}), biasanya akan ada gelas dan teko air putih di sudut ruangan. Air putih gratis ini disebut \textbf{お冷 (O-hiya)}. Silakan ambil sendiri untuk seluruh keluargamu!
\end{lessonbox}

% ================= PAIR 36 =================
\begin{convobox}{71}
    \noindent
    \jpkana{【あなた】コーヒー}{[Anata] Koohii} \jpkana{ひとつ}{hitotsu} \jpkana{と、}{to,} 
    \jpkana{アップルジュース}{appuru juusu} \jpkana{ひとつ}{hitotsu} \jpkana{と}{to} 
    \jpkana{ブルーベリーマフィン}{buruuberii mafin} \jpkana{ふたつ}{futatsu} 
    \jpkana{お}{o}\jpword{ねが}{願}{nega}\jpkana{いします。}{ishimasu.}

    \vspace{1.5em}
    \noindent{\Large \color{articolor} \textbf{Arti:} [Kamu] Tolong 1 kopi, dan 1 jus apel, dan 2 muffin blueberry.}
\end{convobox}

\begin{convobox}{72}
    \noindent
    \jpkana{【店員】コーヒーは}{[Tenin] Koohii wa} 
    \jpkana{ホットか}{hotto ka} \jpkana{アイス、}{aisu,} 
    \jpkana{どちらに}{dochira ni} \jpkana{なさいますか？}{nasaimasu ka?}

    \vspace{1.5em}
    \noindent{\Large \color{articolor} \textbf{Arti:} [Kasir] Untuk kopinya, Hot atau Ice, Anda ingin pilih yang mana?}
\end{convobox}

\begin{lessonbox}{🧁 Memesan Kue yang Menggoda}
    Lihat kan panjangnya pesanan di atas? Kamu menggabungkan Kopi, Jus Apel, dan 2 Muffin sekaligus memakai partikel \textbf{と (To)}. Ini adalah pesanan sempurna untuk Ibu, Anak, dan cemilan bersama!
\end{lessonbox}

% ================= PAIR 37 =================
\begin{convobox}{73}
    \noindent
    \jpkana{【あなた】アイスで}{[Anata] Aisu de} 
    \jpkana{お}{o}\jpword{ねが}{願}{nega}\jpkana{いします。}{ishimasu.}

    \vspace{1.5em}
    \noindent{\Large \color{articolor} \textbf{Arti:} [Kamu] Tolong yang Es.}
\end{convobox}

\begin{convobox}{74}
    \noindent
    \jpkana{【店員】お}{[Tenin] O}\jpword{さとう}{砂糖}{satou} \jpkana{・}{・} \jpkana{ミルクは}{miruku wa} 
    \jpword{りよう}{利用}{riyou} \jpkana{ですか？}{desu ka?}

    \vspace{1.5em}
    \noindent{\Large \color{articolor} \textbf{Arti:} [Kasir] Apakah menggunakan gula dan susu?}
\end{convobox}

\begin{lessonbox}{🧋 Pesan Minuman Kekinian: "Matcha Frapucino"}
    Kalau kamu masuk ke gerai kopi besar, minuman es kocok yang populer bukan cuma disebut 'Aisu'. Mereka sering memakai kata 'Frape' atau nama merek. Anak-anak pasti sangat suka minuman dingin manis ini di musim panas Jepang!
\end{lessonbox}

% ================= PAIR 38 =================
\begin{convobox}{75}
    \noindent
    \jpkana{【あなた】いえ、}{[Anata] Ie,} 
    \jpword{だいじょうぶ}{大丈夫}{daijoubu} \jpkana{です。}{desu.}

    \vspace{1.5em}
    \noindent{\Large \color{articolor} \textbf{Arti:} [Kamu] Tidak, tidak usah.}
\end{convobox}

\begin{convobox}{76}
    \noindent
    \jpkana{【店員】お}{[Tenin] O}\jpword{かいけい}{会計}{kaikei}\jpkana{、}{,} 
    \jpword{にせんえん}{2,000円}{ni-sen-en} 
    \jpkana{でございます。}{de gozaimasu.}

    \vspace{1.5em}
    \noindent{\Large \color{articolor} \textbf{Arti:} [Kasir] Totalnya 2,000 yen.}
\end{convobox}

\begin{lessonbox}{💴 Membayar Tanpa Ribet!}
    Angka 2.000 Yen dibaca \textbf{Ni-sen-en}. Kalau kamu pakai uang kertas pecahan \textbf{5.000 Yen (Go-sen-en)} atau \textbf{10.000 Yen (Ichi-man-en)}, berikan saja ke kasir. Mereka punya mesin penghitung otomatis yang akan mengeluarkan uang kembalianmu dengan sangat akurat! 
\end{lessonbox}

% ================= PAIR 39 =================
\begin{convobox}{77}
    \noindent
    \jpkana{いかがでしたか？}{Ikaga deshita ka?}

    \vspace{1.5em}
    \noindent{\Large \color{articolor} \textbf{Arti:} Bagaimana (pelajaran hari ini)?}
\end{convobox}

\begin{convobox}{78}
    \noindent
    \jpword{こんど}{今度}{Kondo} \jpword{にほん}{日本}{Nihon} \jpkana{で}{de} 
    \jpkana{カフェに}{kafe ni} 
    \jpword{い}{行}{i}\jpkana{くことがあったら、}{ku koto ga attara,}

    \vspace{1.5em}
    \noindent{\Large \color{articolor} \textbf{Arti:} Jika lain kali kamu ada kesempatan pergi ke kafe di Jepang,}
\end{convobox}

\begin{lessonbox}{✈️ Petualangan Menantimu!}
    Jepang memiliki banyak sekali 'Themed Cafe' (Kafe bertema) lucu seperti Pokemon Cafe, Kirby Cafe, atau Sanrio Cafe yang pastiiiiii bikin anak-anak melonjak kegirangan! Jangan lupa reservasi (pesan tempat) dari jauh hari karena kafe tematik selalu penuh!
\end{lessonbox}

% ================= PAIR 40 =================
\begin{convobox}{79}
    \noindent
    \jpkana{ぜひ}{Zehi} 
    \jpkana{この}{kono} \jpword{どうが}{動画}{douga} \jpkana{で}{de} 
    \jpword{まな}{学}{mana}\jpkana{んだ}{nda} 
    \jpkana{フレーズを}{fureezu o} 
    \jpword{つか}{使}{tsuka}\jpkana{って}{tte} 
    \jpword{ちゅうもん}{注文}{chuumon} \jpkana{してみてくださいね。}{shite mite kudasai ne.}

    \vspace{1.5em}
    \noindent{\Large \color{articolor} \textbf{Arti:} Pastikan (jangan ragu) untuk memesan dengan menggunakan frasa-frasa yang telah dipelajari di video ini ya.}
\end{convobox}

\begin{convobox}{80}
    \noindent
    \jpkana{それでは}{Sore dewa} 
    \jpkana{また}{mata} 
    \jpword{じかい}{次回}{jikai} \jpkana{の}{no} 
    \jpword{どうが}{動画}{douga} \jpkana{で}{de} 
    \jpkana{お}{o}\jpword{あ}{会}{a}\jpkana{いしましょう！}{i shimashou!}

    \vspace{1.5em}
    \noindent{\Large \color{articolor} \textbf{Arti:} Kalau begitu, sampai jumpa lagi di video (materi) berikutnya!}
\end{convobox}

% --- SUMMARY BOX ---
\vspace{2em}
\begin{tcolorbox}[
    colback=blue!5!white,
    colframe=blue!80!black,
    boxrule=3pt,
    arc=5mm,
    title={\huge\bfseries \sffamily 🌟 Rangkuman Misi Kafe Keluarga 🌟},
    fonttitle=\color{white},
    attach boxed title to top center={yshift=-4mm},
    boxed title style={colback=blue!80!black, rounded corners},
    left=5mm, right=5mm, top=7mm, bottom=5mm
]
\Large \textbf{Hebat! Keluarga Kalian Sudah Siap Nongkrong di Kafe Tokyo!}
\begin{itemize}
    \item \textbf{Cari Meja Dulu:} Selalu amankan mejamu sebelum maju memesan.
    \item \textbf{Pajaknya Beda:} \textit{Ten-nai} (Makan di tempat) dan \textit{Mochikaeri} (Bawa pulang) memiliki aturan pajak yang berbeda di kasir.
    \item \textbf{Rumus Pesanan:} [Nama Menu] + [Jumlah: Hitotsu/Futatsu] + \textit{Onegaishimasu}.
    \item \textbf{Pertanyaan Suhu:} \textit{Hotto} (Panas) atau \textit{Aisu} (Dingin).
    \item \textbf{Customizing:} \textit{... Dake} (Hanya ...), atau \textit{Ie, daijoubu desu} (Tolak dengan halus).
    \item \textbf{Etika Kebersihan:} Kembalikan nampan dan gelas kotor ke area 'Henkyakuguchi'.
\end{itemize}
Dengan berlatih mengucapkan ini keras-keras di rumah, kalian pasti akan memesan makanan seperti warga lokal Jepang sungguhan. Selamat berlibur dan menikmati kopi serta jus yang lezat!
\end{tcolorbox}

\newpage